\documentclass[a4paper,12pt]{article}
\usepackage{hyperref}
\usepackage{amsmath}
\usepackage{geometry}
\geometry{margin=1in}

\title{\textbf{iCalendar Syntax: A Comprehensive Guide}}
\author{}
\date{}

\begin{document}

\maketitle

\tableofcontents
\newpage

\section{Introduction}
The iCalendar format (RFC 5545) is a widely-used standard for calendar data exchange. It is designed for interoperability among calendar systems, including CalDAV-based platforms like Apple Calendar, Google Calendar, and more. This document provides an in-depth explanation of iCalendar syntax and its components, with examples and best practices.

\section{iCalendar Structure}
An iCalendar file consists of text data encapsulated between \texttt{BEGIN:VCALENDAR} and \texttt{END:VCALENDAR}. It contains metadata about the calendar itself and one or more events, tasks, or other calendar objects.

\subsection{Basic Structure}
\begin{verbatim}
BEGIN:VCALENDAR
VERSION:2.0
PRODID:-//YourApp//Example Calendar 1.0//EN
CALSCALE:GREGORIAN
BEGIN:VEVENT
UID:unique-id@example.com
SUMMARY:Event Title
DTSTART:YYYYMMDDTHHMMSSZ
DTEND:YYYYMMDDTHHMMSSZ
END:VEVENT
END:VCALENDAR
\end{verbatim}

\section{Key Components}
\subsection{VERSION}
Defines the iCalendar version. The current standard is \texttt{2.0}.

\subsection{PRODID}
Specifies the identifier of the software that generated the iCalendar file. Example:
\begin{verbatim}
PRODID:-//YourApp//Example Calendar 1.0//EN
\end{verbatim}

\subsection{CALSCALE}
Defines the calendar scale. The default is \texttt{GREGORIAN}, which represents the Gregorian calendar system. Example:
\begin{verbatim}
CALSCALE:GREGORIAN
\end{verbatim}

\section{Event Definition}
An event is defined within \texttt{BEGIN:VEVENT} and \texttt{END:VEVENT} blocks. Below are the key fields:

\subsection{UID}
A unique identifier for the event. This ensures that the event can be updated or removed across systems. Example:
\begin{verbatim}
UID:unique-id@example.com
\end{verbatim}

\subsection{SUMMARY}
A short title for the event. Example:
\begin{verbatim}
SUMMARY:Weekly Meeting
\end{verbatim}

\subsection{DTSTART and DTEND}
Defines the start and end times of the event. Use the format \texttt{YYYYMMDDTHHMMSSZ} for date-time in UTC or \texttt{YYYYMMDD} for all-day events.

\textbf{Example:} An event starting at 9:00 AM UTC on January 15, 2025:
\begin{verbatim}
DTSTART:20250115T090000Z
DTEND:20250115T100000Z
\end{verbatim}

\subsection{DTSTAMP}
Records the timestamp when the event was created or last modified. Example:
\begin{verbatim}
DTSTAMP:20250114T120000Z
\end{verbatim}

\subsection{SEQUENCE}
Tracks revisions of the event. Increment this number each time the event is modified. Example:
\begin{verbatim}
SEQUENCE:0
\end{verbatim}

\subsection{STATUS}
Defines the state of the event. Options include:
\begin{itemize}
    \item \texttt{TENTATIVE}: The event is not confirmed.
    \item \texttt{CONFIRMED}: The event is finalized.
    \item \texttt{CANCELLED}: The event has been canceled.
\end{itemize}
\textbf{Example:}
\begin{verbatim}
STATUS:CONFIRMED
\end{verbatim}

\subsection{TRANSP}
Indicates whether the event blocks your calendar availability. Options include:
\begin{itemize}
    \item \texttt{OPAQUE}: Blocks availability.
    \item \texttt{TRANSPARENT}: Does not block availability.
\end{itemize}
\textbf{Example:}
\begin{verbatim}
TRANSP:OPAQUE
\end{verbatim}

\subsection{CATEGORIES}
Tags or categorizes the event. Example:
\begin{verbatim}
CATEGORIES:Work,Meeting
\end{verbatim}

\section{Recurrence Rules (RRULE)}
\texttt{RRULE} defines how events repeat. The syntax is:
\begin{verbatim}
RRULE:FREQ=<frequency>;INTERVAL=<interval>;BYDAY=<days>;UNTIL=<end-date>
\end{verbatim}

\subsection{Options}
\begin{itemize}
    \item \texttt{FREQ}: Frequency of recurrence (\texttt{DAILY}, \texttt{WEEKLY}, \texttt{MONTHLY}, \texttt{YEARLY}).
    \item \texttt{INTERVAL}: Number of intervals between occurrences.
    \item \texttt{BYDAY}: Days of the week (\texttt{MO}, \texttt{TU}, \texttt{WE}, etc.).
    \item \texttt{UNTIL}: End date of the recurrence (\texttt{YYYYMMDDTHHMMSSZ} format).
\end{itemize}

\subsection{Example}
A weekly event on Wednesday and Thursday, starting January 15, 2025, and ending February 28, 2025:
\begin{verbatim}
RRULE:FREQ=WEEKLY;INTERVAL=1;BYDAY=WE,TH;UNTIL=20250228T235959Z
\end{verbatim}

\section{Complete Example}
Below is a complete example of an iCalendar file for a recurring event:
\begin{verbatim}
BEGIN:VCALENDAR
VERSION:2.0
PRODID:-//YourApp//Example Calendar 1.0//EN
CALSCALE:GREGORIAN
BEGIN:VEVENT
UID:unique-id-1234@example.com
SEQUENCE:0
STATUS:CONFIRMED
TRANSP:OPAQUE
SUMMARY:Weekly Meeting
DESCRIPTION:Weekly meeting for project updates.
LOCATION:Conference Room A
DTSTART:20250115T090000Z
DTEND:20250115T100000Z
RRULE:FREQ=WEEKLY;INTERVAL=1;BYDAY=WE,TH;UNTIL=20250228T235959Z
DTSTAMP:20250114T120000Z
CATEGORIES:Work,Meeting
END:VEVENT
END:VCALENDAR
\end{verbatim}

\section{Conclusion}
The iCalendar format is a powerful tool for managing calendar data across systems. By understanding its syntax and fields, you can create, modify, and manage events programmatically. This guide provides the foundation to work with iCalendar files effectively, particularly in CalDAV-based environments.

\end{document}